% !TEX root = ../main.tex
\subsection{Design Insight}
Queue length identifies overloaded cores but not which task is worth
rescuing or which target can absorb it safely. RescueSched therefore uses
queue length only as a risk trigger; the final decision depends on the
local counterfactual, the remote counterfactual, migration cost, remote
deadline slack, and incremental target-side risk.

\subsection{Scheduling Loop}
On each scheduling epoch (or queue-update event), RescueSched updates
per-core statistics, identifies risky source cores and low-risk targets,
and for each source scans the first $K$ \emph{locally doomed} candidates.
For each candidate it tests up to $H$ targets, keeping pairs that are
remotely feasible and target-safe, scores them, and commits at most $B$
descriptor-only migrations per epoch, updating simulated target state
between selections.

\begin{algorithm}[t]
\caption{RescueSched main loop (per epoch $t$)}
\label{alg:main}
\begin{algorithmic}[1]
\STATE update $R_c,W_c,Risk(c,t)$ for all $c$
\STATE $S_{risky}\!\gets\!\{s: Risk(s,t)\!>\!0 \lor W_s \text{ over trigger}\}$
\STATE $T_{safe}\!\gets\!\{j: Risk(j,t)\!=\!0 \lor W_j \text{ under trigger}\}$
\STATE $Pairs \gets \emptyset$
\FOR{$s \in S_{risky}$}
  \FOR{$i \in$ top-$K$ locally-doomed of $Q_s$}
    \IF{$\hat{C}_{\text{local}}(i,s)\le D_i$} \STATE \textbf{continue} \ENDIF
    \FOR{$j \in$ top-$H$ of $T_{safe}$, $j\neq s$}
      \IF{$\hat{C}_{\text{remote}}(i,j)+\epsilon > D_i$} \STATE \textbf{continue} \ENDIF
      \IF{$\Delta Risk(j|i,t) > \theta$} \STATE \textbf{continue} \ENDIF
      \STATE $Pairs \mathrel{+}= (i,j,\textsc{Score}(i,j))$
    \ENDFOR
  \ENDFOR
\ENDFOR
\STATE sort $Pairs$ by descending score; $Selected \gets \emptyset$
\FOR{$(i,j)\in Pairs$}
  \IF{$|Selected|\!\ge\!B$} \STATE \textbf{break} \ENDIF
  \IF{$i$ selected $\lor$ $j$ no longer safe} \STATE \textbf{continue} \ENDIF
  \STATE $Selected \mathrel{+}= (i,j)$; update simulated state of $j$
\ENDFOR
\STATE execute descriptor-only migrations for $Selected$
\end{algorithmic}
\end{algorithm}

\subsection{Scoring}
With $\gamma{=}0,\theta{=}0$ in the minimal version, the score reduces to
$Score(i,j)=Slack_{\text{remote}}(i,j)-\beta M_{s,j}(i)$, ranking
migrations by robustness to estimation error net of cost. The full form
adds $\gamma\cdot local\_lateness-\lambda\cdot\Delta Risk$.

\subsection{Complexity}
With per-core prefix sums, $\hat{C}_{\text{local}}$ is $O(1)$ per request
after snapshotting, and target risk is approximated in $O(KH)$ per epoch,
giving naive per-epoch cost $O(|S_{risky}|\cdot K\cdot H\cdot L)$.

\subsection{Migration Usefulness Taxonomy}
Each committed migration is labeled post-hoc as \emph{beneficial},
\emph{needless}, \emph{unsaved}, or \emph{harmful}, yielding the
beneficial-migration ratio (BMR), useless-migration ratio (UMR), and
rescue-per-migration (RPM). These measure whether the scheduler migrates
\emph{precisely}, not merely frequently.